\chapter{The materials}
\label{app:materials}

All ten materials are released under the Creative Commons CC0~1.0 licence (public domain
dedication), which allows any use without attribution; the program and this guide credit them
anyway. Each material is used at 1K resolution ($1024\times1024$).

\begin{table}[ht]
\centering\small
\begin{tabularx}{\linewidth}{@{}llX@{}}
\toprule
\textbf{Material} & \textbf{Source} & \textbf{Page}\\
\midrule
Bricks 076 C      & ambientCG  & \url{https://ambientcg.com/view?id=Bricks076C}\\
Wood 066          & ambientCG  & \url{https://ambientcg.com/view?id=Wood066}\\
Rock 030          & ambientCG  & \url{https://ambientcg.com/view?id=Rock030}\\
Metal 032         & ambientCG  & \url{https://ambientcg.com/view?id=Metal032}\\
Tiles 101         & ambientCG  & \url{https://ambientcg.com/view?id=Tiles101}\\
Fabric 048        & ambientCG  & \url{https://ambientcg.com/view?id=Fabric048}\\
Paving Stones 131 & ambientCG  & \url{https://ambientcg.com/view?id=PavingStones131}\\
Castle Brick 07   & Poly Haven (Rob Tuytel) & \url{https://polyhaven.com/a/castle_brick_07}\\
Aerial Rocks 02   & Poly Haven (Rob Tuytel) & \url{https://polyhaven.com/a/aerial_rocks_02}\\
Bark Willow 02    & Poly Haven (Charlotte Baglioni) & \url{https://polyhaven.com/a/bark_willow_02}\\
\bottomrule
\end{tabularx}
\caption{The materials of Texture Explorer and their origin.}
\end{table}

\chapter{Building this document}
\label{app:build}

The sources of this guide are in the \file{Docs} folder of the repository:
\begin{description}[font=\normalfont\ttfamily]
\item[guide.tex] the main file, with the preamble and the list of chapters;
\item[chapters/] one file per chapter;
\item[figures/] screenshots of the program (the diagrams are drawn with TikZ and pgfplots in
the chapters themselves);
\item[excerpts.txt] the list of code excerpts: for each, a source file and a regular expression
for its first line;
\item[build.ps1] the build script.
\end{description}
The build script first extracts each excerpt from the current sources into \file{build/excerpts},
following the braces of the function or structure that starts at the matching line. It then runs
\code{latexmk}, \code{pdflatex} or \code{tectonic}, whichever it finds, and writes
\file{TextureExplorer-Guide.pdf}:
\begin{lstlisting}[language={},numbers=none,basicstyle=\ttfamily\small]
powershell -ExecutionPolicy Bypass -File Docs\build.ps1
powershell -ExecutionPolicy Bypass -File Docs\build.ps1 -TeX C:\tools\tectonic.exe
\end{lstlisting}
On macOS or Linux the same script runs with PowerShell~7 (\code{brew install --cask powershell},
and \code{brew install tectonic} for a TeX engine):
\begin{lstlisting}[language={},numbers=none,basicstyle=\ttfamily\small]
pwsh Docs/build.ps1
\end{lstlisting}
Because the listings are extracted at build time, the guide always shows the code as it is. If a
function is renamed, the build stops with an error that names the excerpt to fix.

The repository also contains, in \file{LiterateP}, the complete program written as a literate
program in the style of Donald Knuth, from which every source file can be regenerated. That
document follows the code; this one follows the theory.
